% Query execution

\subsection{Query Execution}

Query execution in the sheaf database follows a plan generated by the SheafQueryPlanner, which classifies each query along a local-to-global spectrum and selects the optimal execution strategy.

\subsubsection{Query Classification}

The SheafOptimizer classifies queries into three levels:

\begin{enumerate}
    \item \textbf{Local}: Queries with a specific subject ID (LOOKUP type). The optimizer identifies the entity open set $U_{\text{entity}(s)}$ and retrieves all sections defined on it. This requires no restriction or gluing — it is a direct stalk access.

    \item \textbf{Semi-local}: Queries constrained by context, temporal range, or provenance (CONTEXT, TEMPORAL, PROVENANCE, NEIGHBORHOOD types). The optimizer identifies relevant open sets and chains restriction operations to narrow the result.

    \item \textbf{Global}: Unconstrained queries (GLOBAL, MIXED, AGGREGATION types). The plan either retrieves cached global sections or computes them via the gluing algorithm.
\end{enumerate}

\subsubsection{Execution Plan}

Each plan consists of:

\begin{enumerate}
    \item \textbf{Classify}: Determine the query level and candidate open sets.
    \item \textbf{Retrieve/Restrict}: For local queries, directly retrieve sections. For semi-local queries, apply restriction maps. For global queries, compute or retrieve global sections.
    \item \textbf{Filter and Limit}: Apply any remaining filters (subject, relation, context) and enforce limit/offset.
\end{enumerate}

\subsubsection{Optimization Rules}

The optimizer follows a strict preference hierarchy:

\begin{enumerate}
    \item Prefer local evaluation over semi-local over global.
    \item Cache reusable global sections; invalidate on write.
    \item Avoid repeated restriction computations by tracking restriction application counts.
    \item Deduplicate candidate open sets before execution.
\end{enumerate}

\subsubsection{Complexity Analysis}

\begin{itemize}
    \item \textbf{Local query}: $O(1)$ open set lookup, $O(|F(U)|)$ section retrieval.
    \item \textbf{Semi-local query}: $O(d \cdot |F(U)|)$ where $d$ is the depth of restriction chains.
    \item \textbf{Global query}: $O(k^2 \cdot m)$ with caching (first invocation) or $O(1)$ with cache hit.
\end{itemize}

\noindent This compares favorably to equivalent KG queries which require $O(n)$ triple index lookups and $O(j)$ join operations for $n$-ary facts reconstructed from $j = n-1$ triples.
\endinput
